%%%%%%%%%%%%%%%%%%%%%%%%%%%%%%%%%%%%%%%%%
% Bachelor/Master Thesis 
% LaTeX Template
% Version 2.5 (27/8/17)
%
% This template was downloaded from:
% http://www.LaTeXTemplates.com
%
% Version 2.x major modifications by:
% Vel (vel@latextemplates.com)
%
% This template is based on a template by:
% Steve Gunn (http://users.ecs.soton.ac.uk/srg/softwaretools/document/templates/)
% Sunil Patel (http://www.sunilpatel.co.uk/thesis-template/)
%
% Template license:
% CC BY-NC-SA 3.0 (http://creativecommons.org/licenses/by-nc-sa/3.0/)
%
%%%%%%%%%%%%%%%%%%%%%%%%%%%%%%%%%%%%%%%%%

%----------------------------------------------------------------------------------------
%	PACKAGES AND OTHER DOCUMENT CONFIGURATIONS
%----------------------------------------------------------------------------------------
% https://www.overleaf.com/9572798646vszsfmjxtyqc
\documentclass[
11pt, % The default document font size, options: 10pt, 11pt, 12pt
%oneside, % Two side (alternating margins) for binding by default, uncomment to switch to one side
english, % ngerman for German
singlespacing, % Single line spacing, alternatives: onehalfspacing or doublespacing
%draft, % Uncomment to enable draft mode (no pictures, no links, overfull hboxes indicated)
%nolistspacing, % If the document is onehalfspacing or doublespacing, uncomment this to set spacing in lists to single
%liststotoc, % Uncomment to add the list of figures/tables/etc to the table of contents
%toctotoc, % Uncomment to add the main table of contents to the table of contents
%parskip, % Uncomment to add space between paragraphs
%nohyperref, % Uncomment to not load the hyperref package
headsepline, % Uncomment to get a line under the header
%chapterinoneline, % Uncomment to place the chapter title next to the number on one line
%consistentlayout, % Uncomment to change the layout of the declaration, abstract and acknowledgements pages to match the default layout
oneside,]{BachelorMasterThesis} % The class file specifying the document structure

% \usepackage{graphicx} % for graphics
\usepackage{float} % for graphics if you want the figure in a particular place (add [H])
\usepackage{subcaption} % for sub figures

\usepackage[utf8]{inputenc}
\usepackage{epigraph}

\usepackage[utf8]{inputenc} % Required for inputting international characters
\usepackage[T1]{fontenc} % Output font encoding for international characters

\usepackage{mathpazo} % Use the Palatino font by default

\usepackage[backend=bibtex,style=authoryear,natbib=true]{biblatex} % Use the bibtex backend with the authoryear citation style (which resembles APA)

\usepackage{listings} % for code snippets

\addbibresource{rmod,others} % The filename of the bibliography

\usepackage[autostyle=true]{csquotes} % Required to generate language-dependent quotes in the bibliography

%\usepackage[dvipsnames]{xcolor} %%% for comments
\usepackage{ifthen} %%% for comments
\usepackage{amssymb} %%% for comments
\usepackage{soul} %%% for strikethrough

\usepackage{enumitem}
\newlist{RQ}{enumerate}{1}
\setlist[RQ]{label=RQ\arabic*:}

% comments \nb{label}{color}{text}
\newboolean{showcomments}
\setboolean{showcomments}{true}
\ifthenelse{\boolean{showcomments}}
    {\newcommand{\nb}[3]{
        {\colorbox{#2}{\bfseries\sffamily\scriptsize\textcolor{white}{#1}}}
        {\textcolor{#2}{\sf\small$\blacktriangleright$\textit{#3}$\blacktriangleleft$}}}
     \newcommand{\version}{\emph{\scriptsize$-$Id$-$}}}
    {\newcommand{\nb}[3]{}
     \newcommand{\version}{}}
\newcommand{\oz}[1]{\nb{Oleks}{blue}{#1}}
\newcommand{\md}[1]{\nb{Marcus}{red}{#1}}
\newcommand{\mr}[1]{\nb{Myroslava}{olive}{#1}}

\newcommand{\red}[1]{\textcolor{red}{#1}}
\newcommand{\green}[1]{\textcolor{OliveGreen}{#1}}
\newcommand{\blue}[1]{\textcolor{blue}{#1}}
\newcommand{\yellow}[1]{\textcolor{orange}{#1}}
\newcommand{\purple}[1]{\textcolor{Plum}{#1}}

\newcommand{\notclear}[1]{\yellow{#1}}
\newcommand{\wrong}[1]{\red{#1}}
\newcommand{\remove}[1]{\red{\st{#1}}}
\newcommand{\replace}[2]{\red{\st{#1} [#2]}}
\newcommand{\insertion}[1]{\red{[#1]}}
\newcommand{\badstyle}[1]{\purple{#1}}
\newcommand{\repetition}[1]{\green{#1}}

%%% for smalltalk-syntax highlighting
% Just add % Just add % Just add \input{smalltalkEnv} to your file
% then you can use :
% \begin{lstlisting}[language=Smalltalk]
% false become: true.
% \end{lstlisting}

\usepackage{color}
\usepackage{listings}
\usepackage{etoolbox}
\usepackage{textcomp}

\definecolor{stComment}{rgb}{0.5,0.5,0.5}
\definecolor{stString}{rgb}{0.58,0,0.82}
\definecolor{stKeywords}{rgb}{0.21,0.55,0.7}
\definecolor{stNumbers}{rgb}{.5,0,0}

\newtoggle{InString}{}% Keep track of if we are within a string
\togglefalse{InString}% Assume not initally in string
\newcommand*{\ColorIfNotInString}[1]{\iftoggle{InString}{#1}{\color{stNumbers}#1}}%
\newcommand*{\ProcessQuote}[1]{#1\iftoggle{InString}{\global\togglefalse{InString}}{\global\toggletrue{InString}}}%

\lstdefinelanguage{Smalltalk}{
  keywordstyle=\color{stKeywords},
  commentstyle=\color{stComment},
  stringstyle=\color{stString},
  alsoletter=\#,
  identifierstyle=\idstyle, 
  showstringspaces=false,
  morekeywords={true,false,self,super,nil},
  sensitive=true, 
  morecomment=[s]{"}{"}, 
  morestring=[d]', 
  style=SmalltalkStyle,
  tabsize=2,
  basicstyle=\ttfamily,
  upquote=true,
}


\makeatletter%
\newcommand*\idstyle[1]{%
  \expandafter\id@style\the\lst@token{#1}\relax%
}
\def\id@style#1#2\relax{%
  \ifnum\pdfstrcmp{#1}{\#}=0%
  \ttfamily\color{stString} \the\lst@token%
  \else%
  \edef\tempa{\uccode`#1}%
  \edef\tempb{`#1}%
  \ifnum\tempa=\tempb%
  \ttfamily\color{blue} \the\lst@token%
  \else%
  \the\lst@token%
  \fi%
  \fi%
}

\lstdefinestyle{SmalltalkStyle}{ 
  literate=%
  {^}{{$\uparrow$}}1% 
  % {"}{{{\ProcessQuote{"}}}}1% Disable coloring within double quotes
  % {'}{{{\ProcessQuote{'}}}}1% Disable coloring within single quote
  {0}{{{\ColorIfNotInString{0}}}}1%
  {1}{{{\ColorIfNotInString{1}}}}1%
  {2}{{{\ColorIfNotInString{2}}}}1%
  {3}{{{\ColorIfNotInString{3}}}}1%
  {4}{{{\ColorIfNotInString{4}}}}1%
  {5}{{{\ColorIfNotInString{5}}}}1%
  {6}{{{\ColorIfNotInString{6}}}}1%
  {7}{{{\ColorIfNotInString{7}}}}1%
  {8}{{{\ColorIfNotInString{8}}}}1%
  {9}{{{\ColorIfNotInString{9}}}}1%
}  to your file
% then you can use :
% \begin{lstlisting}[language=Smalltalk]
% false become: true.
% \end{lstlisting}

\usepackage{color}
\usepackage{listings}
\usepackage{etoolbox}
\usepackage{textcomp}

\definecolor{stComment}{rgb}{0.5,0.5,0.5}
\definecolor{stString}{rgb}{0.58,0,0.82}
\definecolor{stKeywords}{rgb}{0.21,0.55,0.7}
\definecolor{stNumbers}{rgb}{.5,0,0}

\newtoggle{InString}{}% Keep track of if we are within a string
\togglefalse{InString}% Assume not initally in string
\newcommand*{\ColorIfNotInString}[1]{\iftoggle{InString}{#1}{\color{stNumbers}#1}}%
\newcommand*{\ProcessQuote}[1]{#1\iftoggle{InString}{\global\togglefalse{InString}}{\global\toggletrue{InString}}}%

\lstdefinelanguage{Smalltalk}{
  keywordstyle=\color{stKeywords},
  commentstyle=\color{stComment},
  stringstyle=\color{stString},
  alsoletter=\#,
  identifierstyle=\idstyle, 
  showstringspaces=false,
  morekeywords={true,false,self,super,nil},
  sensitive=true, 
  morecomment=[s]{"}{"}, 
  morestring=[d]', 
  style=SmalltalkStyle,
  tabsize=2,
  basicstyle=\ttfamily,
  upquote=true,
}


\makeatletter%
\newcommand*\idstyle[1]{%
  \expandafter\id@style\the\lst@token{#1}\relax%
}
\def\id@style#1#2\relax{%
  \ifnum\pdfstrcmp{#1}{\#}=0%
  \ttfamily\color{stString} \the\lst@token%
  \else%
  \edef\tempa{\uccode`#1}%
  \edef\tempb{`#1}%
  \ifnum\tempa=\tempb%
  \ttfamily\color{blue} \the\lst@token%
  \else%
  \the\lst@token%
  \fi%
  \fi%
}

\lstdefinestyle{SmalltalkStyle}{ 
  literate=%
  {^}{{$\uparrow$}}1% 
  % {"}{{{\ProcessQuote{"}}}}1% Disable coloring within double quotes
  % {'}{{{\ProcessQuote{'}}}}1% Disable coloring within single quote
  {0}{{{\ColorIfNotInString{0}}}}1%
  {1}{{{\ColorIfNotInString{1}}}}1%
  {2}{{{\ColorIfNotInString{2}}}}1%
  {3}{{{\ColorIfNotInString{3}}}}1%
  {4}{{{\ColorIfNotInString{4}}}}1%
  {5}{{{\ColorIfNotInString{5}}}}1%
  {6}{{{\ColorIfNotInString{6}}}}1%
  {7}{{{\ColorIfNotInString{7}}}}1%
  {8}{{{\ColorIfNotInString{8}}}}1%
  {9}{{{\ColorIfNotInString{9}}}}1%
}  to your file
% then you can use :
% \begin{lstlisting}[language=Smalltalk]
% false become: true.
% \end{lstlisting}

\usepackage{color}
\usepackage{listings}
\usepackage{etoolbox}
\usepackage{textcomp}

\definecolor{stComment}{rgb}{0.5,0.5,0.5}
\definecolor{stString}{rgb}{0.58,0,0.82}
\definecolor{stKeywords}{rgb}{0.21,0.55,0.7}
\definecolor{stNumbers}{rgb}{.5,0,0}

\newtoggle{InString}{}% Keep track of if we are within a string
\togglefalse{InString}% Assume not initally in string
\newcommand*{\ColorIfNotInString}[1]{\iftoggle{InString}{#1}{\color{stNumbers}#1}}%
\newcommand*{\ProcessQuote}[1]{#1\iftoggle{InString}{\global\togglefalse{InString}}{\global\toggletrue{InString}}}%

\lstdefinelanguage{Smalltalk}{
  keywordstyle=\color{stKeywords},
  commentstyle=\color{stComment},
  stringstyle=\color{stString},
  alsoletter=\#,
  identifierstyle=\idstyle, 
  showstringspaces=false,
  morekeywords={true,false,self,super,nil},
  sensitive=true, 
  morecomment=[s]{"}{"}, 
  morestring=[d]', 
  style=SmalltalkStyle,
  tabsize=2,
  basicstyle=\ttfamily,
  upquote=true,
}


\makeatletter%
\newcommand*\idstyle[1]{%
  \expandafter\id@style\the\lst@token{#1}\relax%
}
\def\id@style#1#2\relax{%
  \ifnum\pdfstrcmp{#1}{\#}=0%
  \ttfamily\color{stString} \the\lst@token%
  \else%
  \edef\tempa{\uccode`#1}%
  \edef\tempb{`#1}%
  \ifnum\tempa=\tempb%
  \ttfamily\color{blue} \the\lst@token%
  \else%
  \the\lst@token%
  \fi%
  \fi%
}

\lstdefinestyle{SmalltalkStyle}{ 
  literate=%
  {^}{{$\uparrow$}}1% 
  % {"}{{{\ProcessQuote{"}}}}1% Disable coloring within double quotes
  % {'}{{{\ProcessQuote{'}}}}1% Disable coloring within single quote
  {0}{{{\ColorIfNotInString{0}}}}1%
  {1}{{{\ColorIfNotInString{1}}}}1%
  {2}{{{\ColorIfNotInString{2}}}}1%
  {3}{{{\ColorIfNotInString{3}}}}1%
  {4}{{{\ColorIfNotInString{4}}}}1%
  {5}{{{\ColorIfNotInString{5}}}}1%
  {6}{{{\ColorIfNotInString{6}}}}1%
  {7}{{{\ColorIfNotInString{7}}}}1%
  {8}{{{\ColorIfNotInString{8}}}}1%
  {9}{{{\ColorIfNotInString{9}}}}1%
} 
\lstset{language=Smalltalk}

%----------------------------------------------------------------------------------------
%	MARGIN SETTINGS
%----------------------------------------------------------------------------------------

\geometry{
	paper=a4paper, % Change to letterpaper for US letter
	inner=2.5cm, % Inner margin
	outer=3.8cm, % Outer margin
	bindingoffset=.5cm, % Binding offset
	top=1.5cm, % Top margin
	bottom=1.5cm, % Bottom margin
	%showframe, % Uncomment to show how the type block is set on the page
}

%----------------------------------------------------------------------------------------
%	THESIS INFORMATION
%----------------------------------------------------------------------------------------

\thesistitle{Improving Code Completion in Pharo Using N-gram Language Models} % Your thesis title, this is used in the title and abstract, print it elsewhere with \ttitle
\supervisor{\href{https://www.oleks.fr/}{Oleksandr \textsc{Zaitsev}} \\ \href{https://marcusdenker.de/}{Dr. Marcus \textsc{Denker}}} % Your supervisor's name, this is used in the title page, print it elsewhere with \supname
\examiner{} % Your examiner's name, this is not currently used anywhere in the template, print it elsewhere with \examname
\degree{Bachelor of Science} % Your degree name, this is used in the title page and abstract, print it elsewhere with \degreename
\author{Myroslava \textsc{Romaniuk}} % Your name, this is used in the title page and abstract, print it elsewhere with \authorname
\addresses{} % Your address, this is not currently used anywhere in the template, print it elsewhere with \addressname

\subject{Software Engineering} % Your subject area, this is not currently used anywhere in the template, print it elsewhere with \subjectname
\keywords{Code Completion} % Keywords for your thesis, this is not currently used anywhere in the template, print it elsewhere with \keywordnames
\university{\href{http://www.ucu.edu.ua}{Ukrainian Catholic University}} % Your university's name and URL, this is used in the title page and abstract, print it elsewhere with \univname
\department{\href{https://apps.ucu.edu.ua/en/}{Faculty of Applied Sciences}} % Your department's name and URL, this is used in the title page and abstract, print it elsewhere with \deptname
\group{Department of Computer Sciences} % Your research group's name and URL, this is used in the title page, print it elsewhere with \groupname
\faculty{\href{http://faculty.university.com}{}} % Your faculty's name and URL, this is used in the title page and abstract, print it elsewhere with \facname

\AtBeginDocument{
\hypersetup{pdftitle=\ttitle} % Set the PDF's title to your title
\hypersetup{pdfauthor=\authorname} % Set the PDF's author to your name
\hypersetup{pdfkeywords=\keywordnames} % Set the PDF's keywords to your keywords
}

\begin{document}

\frontmatter % Use roman page numbering style (i, ii, iii, iv...) for the pre-content pages

\pagestyle{plain} % Default to the plain heading style until the thesis style is called for the body content

%----------------------------------------------------------------------------------------
%	TITLE PAGE
%----------------------------------------------------------------------------------------

\begin{titlepage}
\begin{center}

\vspace*{.06\textheight}
{\scshape\LARGE \univname\par}\vspace{1.5cm} % University name
\textsc{\Large Bachelor Thesis}\\[0.5cm] % Thesis type

\HRule \\[0.4cm] % Horizontal line
{\huge \bfseries \ttitle\par}\vspace{0.4cm} % Thesis title
\HRule \\[1.5cm] % Horizontal line
 
\begin{minipage}[t]{0.4\textwidth}
\begin{flushleft} \large
\emph{Author:}\\
\authorname % Author name - remove the \href bracket to remove the link
\end{flushleft}
\end{minipage}
\begin{minipage}[t]{0.4\textwidth}
\begin{flushright} \large
\emph{Supervisors:} \\
\supname % Supervisor name - remove the \href bracket to remove the link  
\end{flushright}
\end{minipage}\\[3cm]
 
\vfill

\large \textit{The thesis submitted in fulfillment of the requirements\\ for the degree of \degreename}\\[0.3cm] % University requirement text
\textit{in the}\\[0.4cm]
\groupname\\\deptname\\[1cm] % Research group name and department name
 
\vfill
\includegraphics[height=5cm]{UCU-Apps.png} % University/department logo - uncomment to place it

\vfill
{\large Lviv 2020}\\[3cm] % Date
 
\vfill
\end{center}
\end{titlepage}

%----------------------------------------------------------------------------------------
%	DECLARATION PAGE
%----------------------------------------------------------------------------------------

\begin{declaration}
\addchaptertocentry{\authorshipname} % Add the declaration to the table of contents
\noindent I, \authorname, declare that this thesis titled, \enquote{\ttitle} and the work presented in it are my own. I confirm that:

\begin{itemize} 
\item This work was done wholly or mainly while in candidature for a research degree at this University.
\item Where any part of this thesis has previously been submitted for a degree or any other qualification at this University or any other institution, this has been clearly stated.
\item Where I have consulted the published work of others, this is always clearly attributed.
\item Where I have quoted from the work of others, the source is always given. With the exception of such quotations, this thesis is entirely my own work.
\item I have acknowledged all main sources of help.
\item Where the thesis is based on work done by myself jointly with others, I have made clear exactly what was done by others and what I have contributed myself.\\
\end{itemize}
 
\noindent Signed:\\
\rule[0.5em]{25em}{0.5pt} % This prints a line for the signature
 
\noindent Date:\\
\rule[0.5em]{25em}{0.5pt} % This prints a line to write the date
\end{declaration}

\cleardoublepage

%----------------------------------------------------------------------------------------
%	QUOTATION PAGE
%----------------------------------------------------------------------------------------

\vspace*{0.2\textheight}

\noindent\enquote{Please do not be cynical. I hate cynicism. For the record, it's my least favorite quality. It doesn't lead anywhere. No one in life gets exactly what they thought they were going to get. But if you work really hard and you're kind, amazing things will happen.}\bigbreak

\hfill Conan O'Brien

%----------------------------------------------------------------------------------------
%	ABSTRACT PAGE
%----------------------------------------------------------------------------------------

\begin{abstract}
\addchaptertocentry{\abstractname} % Add the abstract to the table of contents
Code completion is one of the essential features of any IDE and it significantly improves the developer experience and productivity. Good code completion can both speed up the development process, as well as aid the developer in API exploration. On the other hand, having a slow or inaccurate completion can be very cumbersome. Thus, it is important to find a way to make it as effective as possible.

The current implementation of code completion in the Pharo IDE is based on the abstract syntax tree (AST) of source code. The AST allows us to learn the semantic role and the kind of tokens (e.g. class name, method name, literal, et cetera) and to suggest contextually relevant completions. However, the current implementation has no efficient way to sort the completion candidates, which means that sometimes the user must scroll through a long list of proposed completions to find the one that they need.

In this work, I study how code completion can be improved by sorting completions in such a way that most frequently used tokens appear first. To that end, I implement a sorting plugin based on the n-gram language model.

I use quantitative and qualitative evaluation techniques to compare the performance of the n-gram (unigram and bigram) and alphabetic sorters. As a result of the evaluation, the unigram sorter was shown to have the best performance.

The sorting strategies implemented are part of a Pharo IDE plugin. It is open source and is available at \url{https://github.com/myroslavarm/CompletionSorting}.
\end{abstract}

%----------------------------------------------------------------------------------------
%	ACKNOWLEDGEMENTS
%----------------------------------------------------------------------------------------

\begin{acknowledgements}
\addchaptertocentry{\acknowledgementname} % Add the acknowledgements to the table of contents
I want to express my endless gratitude to Oleksandr Zaitsev and Dr. Marcus Denker from the RMoD team at INRIA Lille for supervising this work, providing me with valuable feedback, and patiently answering my questions. I also want to thank Dr. Stéphane Ducasse for helping me arrange my internship at INRIA and my stay in Lille, as well as introducing me to Pharo during my first year of university and motivating me to participate in open source development. I am also grateful to INRIA Lille and the European Smalltalk User Group (ESUG) for financing this internship.

I would also like to thank Dr. Yuriy Tymchuk for helping me make my first steps in Pharo during Google Summer of Code (GSoC) 2017, Julien Delplanque and Dr. Guillermo Polito from the RMoD team at INRIA Lille for giving me advice and assisting me throughout various stages of my work on code completion, and Dr. Oles Hodych for providing me with recommendations for writing a thesis and with various \LaTeX\  tips. I am also grateful to my mum and dad for supporting me throughout my Bachelor journey, to my best friends Martin and Ryan for motivating me and keeping me grounded, and to Conan O'Brien and Tina Fey for providing endless comedic relief when it was most needed -- such as writing this thesis in social isolation during the COVID-19 quarantine in a student room in Lille, France.

Finally, I want to thank the Ukrainian Catholic University and the Faculty of Applied Sciences for their continued support throughout my studies and this internship, particularly Oles Dobosevych and Dr. Yaroslav Prytula for their patience, kindness, and belief in me. It has really been an honour to be a part of this program.
\end{acknowledgements}

%----------------------------------------------------------------------------------------
%	LIST OF CONTENTS/FIGURES/TABLES PAGES
%----------------------------------------------------------------------------------------

\tableofcontents % Prints the main table of contents

%----------------------------------------------------------------------------------------
%	ABBREVIATIONS
%----------------------------------------------------------------------------------------

\begin{abbreviations}{ll} % Include a list of abbreviations (a table of two columns)

\textbf{API} & \textbf{A}pplication \textbf{P}rogramming \textbf{I}nterface\\
\textbf{AST} & \textbf{A}bstract \textbf{S}yntax \textbf{T}ree\\
\textbf{IDE} & \textbf{I}ntegrated \textbf{D}evelopment \textbf{E}nvironment\\
\textbf{LSTM} & \textbf{L}ong \textbf{S}hort-\textbf{T}erm \textbf{M}emory\\
\textbf{ML} & \textbf{M}achine \textbf{L}earning\\
\textbf{NLP} & \textbf{N}atural \textbf{L}anguage \textbf{P}rocessing\\
\textbf{RNN} & \textbf{R}ecurrent \textbf{N}eural \textbf{N}etwork\\
\textbf{RQ} & \textbf{R}esearch \textbf{Q}uestion\\

\end{abbreviations}

%----------------------------------------------------------------------------------------
%	THESIS CONTENT - CHAPTERS
%----------------------------------------------------------------------------------------

\mainmatter % Begin numeric (1,2,3...) page numbering

\pagestyle{thesis} % Return the page headers back to the "thesis" style



% Little hack to fix minitoc
%\adjustmtc
%\adjustmtc
%\adjustmtc

% Include the chapters of the thesis as separate files from the Chapters folder
% Uncomment the lines as you write the chapters

\chapter{Introduction}
\label{chap:Introduction}

\section{Context}
\label{sec:Introduction-Context}
Code completion is one of the most used features in any IDE. Whether it is being used for improving speed and accuracy of typing, or is used as an API guide, helping developers find their way around libraries, code completion is an essential part of any successful IDE or a code editor.

Code completion is used as a major productivity tool, and is one of the first things a developer notices as soon as they start coding. The speed with which the results are suggested, as well as their accuracy is very important, and it is something that can certainly "make or break" the workflow of the developer. Therefore, researches and software engineers are always trying to find new approaches that can improv code completion quality and make it effective, as accurate, and as fast \footnote{By "fast" we mean that the completion is quick to give suggestions, by "accurate" that the completion proposes contextually correct results, and by "effective" we mean that the combination of both the aforementioned qualities makes the tool successful in achieving the result, which is to find the most suitable completion as easily as possible.} as possible.

The main distinction among the many approaches to completion code would be whether the approach uses semantic context or not. Semantic models use code structure analysis to tailor completions to type and/or semantic role. Statistical or machine learning models can also often employ semantic analysis but would normally be used solely with code history analysis. In this case, the source code would be treated more like simple text, with semantic roles not taken into consideration.

Moving onto the one code completion engine we are particularly interested in, in the last couple of years there have been several approaches taken to improve the completion inside the Pharo IDE \footnote{The name Pharo refers both to the dynamically typed programming language that is a dialect of Smalltalk, as well as the interactive IDE where this language is used. \newline From now on when referring to the language itself we will simply say \textit{Pharo}, whereas when referring to the IDE we will say \textit{the Pharo IDE} (\url{http://pharo.org/})}.

The current code completion approach is a semantic one: it is based on analysing abstract syntax tree (AST) information of source code, taking the best suitable node, and suggesting specifically tailored completions for each one (i.e. class names for a global node, method calls for a method node, and so on). This is something that we implemented a year ago to replace the less accurate "old" completion engine \mr{need to write more details about how exactly it functioned}. Recently, a separate code completion engine that uses generators was ported to Pharo as well, as a possible alternative version that can be used (the idea is to propose code completion suggestions from pre-generated local contexts using heuristics for optimisation) \footnote{GitHub repository for the project: \url{https://github.com/guillep/complishon}}.

\section{Problem}
\label{sec:Introduction-Problem}
As mentioned above, the current main implementation of code completion in the Pharo IDE is based on parsing and analysing AST information of source code, where we are able to determine the structure of the code, get the semantic role of tokens, as well as infer the type information where possible, with the objective being to only suggest completions that are correct in the given context.

Then, once we have a list of contextually suitable completion candidates, they are sorted alphabetically \footnote{Implementing AST-based completion in Pharo was the focus of an internship project, and alphabetic sorting was made the default through lack of time to implement a more advanced sorting strategy. However, the possibilities of further extending the tool with other sorting strategies were kept in mind, and this bachelor thesis project serves as a continuation of that task.}. This is not the most effective approach, as the process of finding the desired completion candidate can be a bit cumbersome: the correct suggestions can be buried much lower than expected, which would prompt the developer to either scroll through the list, or to type even more characters, taking up valuable time and effort. However, as been mentioned above, the attemps to make code completion more effective are never-ending, and we, too, believe that there are more potential improvements that can be made in the Pharo IDE. In particular, we think the quality of performance of the current completion engine can be greatly improved by taking into account the relevance of suggestions by analysing their use throughout source code history and sorting the results based on that.

\section{Our Approach in a Nutshell}
\label{sec:Introduction-Approach}
Hence, in this work we study how code completion can be improved using statistical language models for sorting the results, on top of an already existing semantic completion. In particular, we use the N-gram language model, as it has been documented to be used for such a task (\cite{Hind12a}), to find the most relevant completion suggestions based on their occurence in code history. We then implement a sorting strategy based on the counted above relevance of suggestions in the Pharo IDE itself. Finally, we evaluate whether this approach indeed enhances the quality of code completion. As a result, we answer the following questions:
\begin{itemize}
    \item Can we improve code completion in the Pharo IDE by sorting candidate completions with an N-gram language model?
    \item Can we build a tool based on a trained N-gram model that would propose completion fast enough to be actually suitable for developer use? \oz{Do you answer this question in your work?}
    \item How can we best numerically evaluate the results of code completion produced by different completion strategies?
\end{itemize}

\section{Contributions}
\label{sec:Introduction-Contributions}
The contributions of this work are as follows:
\begin{itemize}
    \item We develop and train N-gram (in particular, uni- (N=1) and bigram (N=2)) language models on the tokenised source code of a dynamically-typed programming language (Pharo)
    \item We implement the sorting functionality based on the results of the above-mentioned training in an open source IDE
    \item We test whether the combination of semantic-based completion and statistical-based sorting can indeed improve the performance of a code completion engine in an IDE and be effective enough to be used by real developers
    \item To our knowledge, this exact approach of separately considering these two parts (i.e. using one known completion approach for the completion engine and another for the sorting of candidates, and separately assembling them together) has not been done before, and in addition nothing of the sort has been implemented in other Smalltalk-centric IDEs
\end{itemize}

\section{Structure of the Thesis}
\label{sec:Introduction-Structure}
{\hypersetup{linkcolor=black}
\begin{description}
	\item [Chapter \ref{chap:Related Works}. \nameref{chap:Related Works}] \hfill \\
	Here we give a detailed overview of the research done around code completion in the last decade and a half. It includes both the solely software engineering approaches predating the use of machine learning for this task, as well as exclusively machine learning experiments of the recent years.
	\item [Chapter \ref{chap:Completion in Pharo}. \nameref{chap:Completion in Pharo}] \hfill \\
	In this chapter, we give a detailed overview of the implementation details behind the current completion engine in the Pharo IDE, challenges of code completion for a dynamically typed language, as well as describe the idea behind the sorter plugin.
	\item [Chapter \ref{chap:N-gram Background}. \nameref{chap:N-gram Background}] \hfill \\
    Here we provide a theoretical background needed to understand this work, which includes the introduction to the concept of language models in general and N-grams in particular, as well as examples and associated challenges.
    \item [Chapter \ref{chap:Proposed solution}. \nameref{chap:Proposed solution}] \hfill \\
	In this chapter, we take a more detailed look at the idea of our solution, as well as give an extended description of the experiment.
	\item [Chapter \ref{chap:Evaluation}. \nameref{chap:Evaluation}] \hfill \\
	In this chapter, we go into depth about evaluation techniques and challenges, as well as detail the evaluation approach taken and the results of this work.
\end{description}
}
\chapter{Related Works}
\label{chap:RelatedWorks}

\section{Earlier Code Completion Systems}
\label{sec:RelatedWorks-EarlierSystems}
Standard code completions in IDEs used to only rely on language-specific pattern matching, i.e. sorting completions alphabetically based on the symbols already typed in. \cite{Robb08a}, however, showed that code completion can be improved by using program history (program modifications over time). They managed to get good results by prioritising suggestions from recently modified method bodies, and even better results by using per-session vocabulary (changes of the last hour) and merging it with type-based completion.

According to \cite{Bruc09a}, up until 10 years ago code completion systems were exclusively based on type information, with no contextual analysis. The authors countered that by implementing intelligent code completions that learned from examples and had a significantly better performance in terms of relevance of suggestions than other (at the time) common implementations. Their solutions included the frequency based code completion (frequency of use of code), an association rule based completion, and the Best Matching Neighbours code completion (method calls of the closest source snippet found, using a modified k-nearest neighbours algorithm), which was the main contribution of their paper. The BMN based implementation was integrated into Eclipse and demonstrated promising results. It was later extended by \cite{Prok15a} (see \ref{sec:RelatedWorks-DeepLearning}).

\section{Software Naturalness}
\label{sec:RelatedWorks-SoftwareNaturalness}
In the paper "On The Naturalness of Software", \cite{Hind12a} compared source code to natural languages. They claim that code is even more repetitive, predictable and full of patterns than human languages. In the paper, they also argue that code can be modelled by statistical language models, which can be used to support software engineers. Their approach was based on capturing high-level statistical regularity at the n-gram level by taking n-1 previous tokens that are already entered into the text buffer, and attempting to guess the next token. Using this model, it is possible to estimate the most probable sequences of tokens and suggest the most relevant code completions to developers.

This work served as a catalyst for the following research using natural language processing (NLP) for source code. For instance, \cite{Tu14a} also learnt that code "has a high degree of localness, where identifiers (e.g. variable names) are repeated often within close distance" (\cite{Alla18a}). Thus, they applied a cache mechanism that assigns higher probability to tokens that have been observed most recently, and improved the results even further.

\cite{Nguy13a} enhanced the state-of-the-art n-gram approach by incorporating semantic information into code tokens, rather than treating them as text -- i.e. annotating each token with its data type and semantic role if available, which allowed them to increase predictability even further.

\section{Deep Learning for Code Completion}
\label{sec:RelatedWorks-DeepLearning}
In the more recent years, researchers started applying deep learning models such as deep recurrent neural networks (RNN).

For instance, \cite{Hell19a} recorded the results of a case study of 15,000 completions (completion events) for VisualStudio. One of the conclustions they reached is that even though RNNs often outperform n-gram models in typical natural language settings, n-gram models are sometimes a better choice for modeling source code. For example, the deep learner is better at core method invocations but loses on third-party library calls, whereas the n-gram model naturally outperforms it on internal API calls but loses out on the other categories.

\cite{Prok15a} worked on an extensible inference engine for intelligent code completion systems, called PatternBased Bayesian Network (PBN). Eclipse Code Recommenders project adapted the PBN approach for their intelligent call completion. They also tested (evaluating quality, speed and model size) Best Matching Neighbour algorithm, using additional context information for more precise recommendations, and applying clustering techniques to improve model sizes. They conclude that showing the developer hundreds of recommendations may be as ineffective as showing none, and intelligent code completions better target the needs of developers that are unfamiliar with an API.

\cite{Hell17a} introduced a dynamically updatable N-gram model which outperformed both the traditional N-gram models and the deep learning RNN and long short-term memory (LSTM) models.

\cite{Li17a} developed an attention mechanism which exploits the parent-children information of the AST of source code. As correctly predicting out-of-vocabulary (OoV) values in code completion is largely unsuccessful, the authors implemented a pointer mixture network which either generates a new value through an RNN component, or copies an OoV value from local context through a pointer component.

\cite{Rayc14a} implemented a tool called SLANG, which first extracts abstract histories from the data. Then, these histories are fed to a language model such as an n-gram model or recurrent neural network model, which treats the histories as sentences in a natural language and learns probabilities for each sentence, without taking into consideration the AST of the code.

\section{Existing Code Completion Systems in IDEs}
\label{sec:RelatedWorks-ExistingSystems}
The code completion system (Pythia) described by \cite{Svya19a} is part of the Intellicode extension used in the Visual Studio Code IDE to complete Python code. Pythia uses LSTM networks trained on long-range code contexts extracted from abstract syntax trees, which allows it to capture semantics carried by distant nodes and helps rank the method and API recommendations for developers more successfully. Python is a dynamically typed language, so to use type information in Pythia they infer types at runtime based on static analysis of user patterns and add this information into the training sequence.

\cite{Asad14a} developed a tool called CSCC (context-sensitive code completion) is an example-based completion tool which leverages contextual information to better support method call completion. CSCC uses tokenisation instead of deep parsing to collect method call usage patterns and requires type information of the receiver object. The CSCC tool is available as an Eclipse plugin for the Java Editor. Execution time wise the tool showed the results roughly equal to the state-of-the-art, but in terms of method call recommendation accuracy it outperformed those approaches, including BMN (the Best Matching Neighbours implementation from \cite{Bruc09a}).

CACHECA is a cache language model-based code completion tool for Eclipse's Java editor which was presented by \cite{Fran15a}; the tool is based on the cache language model described by \cite{Tu14a} and mentioned earlier in this chapter. According to \cite{Fran15a}, CACHECA greatly enhances Eclipse's built-in engine by incorporating both the corpus and locality statistics, especially when no type information is available.

\section{Summary}
\label{sec:RelatedWorks-Summary}
Here is a brief overview of the progress made in the area of code completion research and development in recent years:

\begin{itemize}
    \item Code completion tools used to only rely on language-specific pattern matching (i.e. by prefix), but around 10 years ago the idea of "intelligent" code completion that can learn from examples was popularised.
    \item The software engineering approaches following this idea relied on frequency of use of code, contextual association, and hierarchical proximity.
    \item From another perspective, source code was likened to natural languages because of its repetitiveness and patterns; the idea that source code can be modelled by statistical language models attracted more machine learning approaches for the task of improving code completion.
    \item Following that, a lot of code completion implementations using machine learning ignored such information as type and semantic meaning, and treated it as regular text, trying to infer relevant predictions from source code history alone.
    \item These days, code completion systems that are actually being used by developers (such as plugins for IDEs) mostly combine both the software engineering approach (taking advantage of semantic and contextual information) and the machine learning one (taking advantage of source code patterns).
\end{itemize}

Another thing worth noting is that the majority of the research experiments mentioned in this chapter have been tested on or applied for only statistically typed languages, predominantly Java and C\# (\cite{Hind12a} also trained on C).

The only related works in this chapter where the experiments included dynamically typed languages were: \cite{Robb08a} (Smalltalk, in addition to Java), \cite{Tu14a} (Python, in addition to Java), \cite{Li17a} (Javascript and Python), and \cite{Svya19a} (Python).
\chapter{Completion in Pharo}
\label{chap:Completion in Pharo}

\section{Code Completion Background}
overview of how code completion works in general and specifically how does it function in the Pharo IDE (challenges, i.e. dynamic language etc)

\section{Existing Pharo Completion}
description of the completion we had pre-AST and what was wrong with it, and how/why we decided to go to AST completion but it's just alphabetical

\section{Sorter Plugin}
explain the role the sorter plays and that it's exactly the part of the existing completion that allows us to experiment with n-gram sorting and so on
\chapter{N-gram Background}
\label{chap:NgramBackground}

In this chapter we go over the theoretical background for the N-grams models, the challenges in their successful implementation, as well as how they can be evaluated.

\section{N-gram Language Models}
\label{sec:NgramBackground-LanguageModels}
In Natural Language Processing (NLP) being able to predict upcoming words or assign probabilities to sequences of words, like phrases and sentences, is one of the essential things in at least several areas, such as speech recognition, spelling correction, spam filtering, machine translation and so on. Thus, by assigning probabilities to words or sequences of words you can improve your accuracy, by suggesting things that are more likely to appear in a given context.

Models that assign such probabilities to sentences or sequences of words are called language models (\cite{Jura09a}). Among them, is the n-gram language model, which assigns probabilities to sequences out of \textit{n} words, called the n-grams. Just one word will be a unigram, 2-gram are referred to as bigrams, 3-grams are referred to as trigrams, and so on. An example of a bigram might be "he ate", "ate the", "the whole" and "whole pizza", whereas a trigram would look like "he ate the", "ate the whole", and "the whole pizza". 

The idea of an n-gram language model could therefore be phrased as the following: to estimate the likelihood of a word in a sequence, based on the "history", which is the previous \textit{(n-1)} words. From this point, we will use the term n-gram to mean either the n-gram language model or the sequence of n words such model uses.

\section{N-grams}
\label{sec:NgramBackground-Ngrams}
To compute the probability of a word \textit{w} based on history \textit{h}, we use the following formula:
\begin{equation}
    P(w|h)
\end{equation}

To put it into a more concrete example, if we have a sentence "he ate the whole pizza", and want to compute the probability of the word "pizza" given that the previous words are "he ate the whole", we can simply do:
\begin{equation}
    P(\text{pizza}|\text{he ate the whole})
\end{equation}

In order to estimate it, we need to compute the number of occurences of "he ate the whole", as well the number occurences of the sentence "he ate the whole pizza", and divide the latter by the former:
\begin{equation}
    P(\text{pizza}|\text{he ate the whole})=\frac{C(\text{he ate the whole pizza})}{C(\text{he ate the whole})}
\end{equation}

However, in large enough data corpora, it can be a challenge to estimate such large sequences. Hence, the idea behind the n-gram model, is that we don't need a large history, but can approximate it using just the last few words (\cite{Jura09a}). This can be done as per the Markov assumption, which means we can reduce our conditional probabilities to be approximately equal:
\begin{equation}
    P(\text{pizza}|\text{he ate the whole}) \sim P(\text{pizza}|\text{whole})
\end{equation}

In general, n-gram probabilities can be estimated by counting the occurences in the corpus and normalising them (as per the maximum likelihood estimate, or MLE for short).

need help with adding MLE and relative frequence explanations?

\section{Evaluating Models}
\label{sec:NgramBackground-EvaluatingModels}
Now, evaluation is a really important step to take after developing and testing a model. One wants to be able to measure the performance and be able to say if it is an improvement on the current state, and by how much. Essentially, there are two ways of evaluation: extrinsic and intrinsic. Extrinsic means actually using the results of the model in an application and measuring how it works. The intrinsic evaluation, however, can measure the results without testing them out in the application, but simply by running the model on the so called test data. It is important to note, that the test data shouldn't be the same as the data the model was initially trained on, to not skew the accuracy measurements.

As raw probability is not used for evaluation language models, when using the intrinsic evaluation something called perplexity is used. Perplexity is the inverse probability of a test set, normalised by the total number of words.

\section{Smoothing}
\label{sec:NgramBackground-Smoothing}
There is, however, the problem of incorrectly estimating probabilities of sequences, if they are not known to us (i.e. weren't in the training set). Those might be perfectly valid sequences that will inadvertently have the probability of 0. Not only can this hurt the performance of an application, but it can also lead to division by 0.

So if the word is entirely unknown, we can convert it to the unknown word token <UNK> and count it as any other word in the dataset. However, if the word we are dealing with is generally known (does exist in the vocabulary) but appears in an unknown context, i.e. after or before a word they never appeared before, we can use a procedure called smoothing, which means we redistibute the general probability a little bit, to not have any sequences with probability equal to 0.

There are several smoothing algorithms, the simplest of which is Laplace smoothing, which we will describe as an example. In this algorithm, one simply needs to add 1 to all the counts, before normalising them (so all the counts will be increased by one, and there will be no 0 counts, and no 0 probabilities). Laplace smoothing doesn't perform well enough to be used in practice. A more useful algorithm based on a similar idea, would be add-k smoothing, which adds a fractional count k (for example, 0.1, 0.5, etc.) to the count, instead of 1.

\section{Summary}
\label{sec:NgramBackground-Summary}
In short:
\begin{itemize}
    \item language models assign probabilities to sentences and sequences of words, and thus predict more likely words
    \item n-gram language models assign probabilities to sequences of \textit{n} words
    \item being a Markov model, the probability of an n-gram model can be approximated to depend just on the several last words, instead of the whole history
    \item n-gram models can be evaluated either extrinsically, by being embedded and tested out in an application, or intrinsically, by testing them on test data and using perplexity
    \item smoothing is needed in case words appear in an unknown context; that way probability is redistributed among all the sequences to prevent skewing of data
\end{itemize}
\chapter{Proposed Solution}
\label{chap:ProposedSolution}

\section{Background}
\label{sec:ProposedSolution-Background}
In our case, in code completion, we want to suggest the most likely tokens at every step. The main idea is, however, that we leave the existing, AST-based code completion engine in place, and just enhance the sorting strategy. That is, after we already have a list of completion candidates proposed to us, we sort it based on each ngram's probability and therefore show the most relevant completions first. In order to do that, we figured it would make sense to train our model on source code data, write it to file, and at every point fetch the relevant, pre-processed information from file where it is stored.

The source code history data used for this was retrieved by \cite{Zait20a} for their research regarding characterising Pharo code. Specifically, the data comes from 50 projects, which consisted of 824 packages, 13,935 classes, and 151,717 methods. In particular, we used the dataset where the source code was already split into tokens and respective token types for each method. In this dataset, even delimiters and non-alphabetic tokens were included, and comments were considered tokens. For example, tokens could include '"Here is a comment for this method"' and '.', and their respective token types would be 'COM' and 'DOT'.

Due to a large number of tokens in our dataset, we needed to be mindful of potential time constraints, as the lookup of ngram probabilities used for sorting had to be fast enough to not make the developer pause and wait for it. Throughout the course of the experiment, we came up with various ways to additionally reduce the dataset, and each of those attempts will be described in the following sections.

\section{Solution Overview}
\label{sec:ProposedSolution-Overview}
We implemented two separate sorting strategies based on the N-gram approach.

The first one, unigram sorting, is based on 1-gram analysis, which means that only the token we are trying to complete is being taken into consideration. The implementation of the unigram model itself comes down to calculating individual token frequencies, i.e. the number of occurrences of each token in the source code history. Then the sorting strategy is implemented by utilising the frequency information and sorting the completion candidates we are given based on each one's number of occurences.

The second, more advanced one, is the bigram sorting strategy. Being an N-gram model with N=2, it relies on both the token currently being completed, as well as the previous history -- in this case, one token before the current one. To make it work, we learnt how to infer the previous token from the context of the method in which we are typing, as well as calculated the probability of each of the ngram sequences. By which we mean that we assembled the history word and each completion candidate into respective pair of tokens and calculated the relative frequency of each one appearing in the source code history, and then based on that returned the probability. The sorting of results per the bigram sorting strategy is thus based on this probability of the generated bigram sequences but is done only for the completion candidates (i.e. words we are predicting in the bigram model), as they are the ones we are interested in seeing displayed in the most relevant order.

The unigram-based sorter is available in the Pharo IDE as the FrequencyCompletionSorter, and the bigram-based as the BigramCompltionSorter -- i.e. they can already be used by enabling the respective sorting option in the Settings. \mr{add screenshot}

\section{Implementation Details}
\label{sec:ProposedSolution-Implementation}
Before training our models, we took some data preparation steps, additionally cleaning the data, such as deleting some rows with token and token type mismatch, eliminating double tabs and replacing them with placeholders, and splitting token and token types into separate columns. After that, when recording token frequencies for the unigram model, we decided we would write the results to file, for faster lookup for the sorting. However, even with that approach, we faced a problem: as the initial number of tokens was around 150k, the lookup was slow enough to be noticeable and break the developer's flow when typing. 

Our approach to improve the time efficiency was to set a certain threshold for the number of occurences, and to cut off all the tokens with frequencies below it. We put the threshold equal to 10, meaning if the token occured less than 10 times throughout the whole history, it was irrelevant enough to be discarded. This way, we were able to cut off most of the miscellaneous, rarely encountered tokens, and shortened our dataset from 150k to just 16k entries. This made the frequencies lookup during completion sorting very fast and, as a result, it was now possible to type and use completion information without any delays.

The one-time tokens, such as custom strings and comments, were not as relevant to us during implementing the unigram sorting strategy, however when we moved on to implementing the bigram sorter, the bigram model became "too heavy" due to many combinations of such tokens, and needed to be additionally reduced. Our solution was to go back to the data processing step and replace those "one-time" tokens with placeholders. In particular, we added placeholders for strings, comments, symbols, characters, arrays and numbers (such as all strings being written as <str> and all numbers as <num>). This allowed us to greatly reduce both the number of total tokens and subsequent ngram sequences, which helped both the bigram and the unigram model, as well as the lookup speed.

\section{Implementation Details: Library Extension}
\label{sec:ProposedSolution-LibraryExtension}
The bigram model trained on source code data was implemented with the help of the NgramModel library \footnote{\url{https://github.com/pharo-ai/NgramModel}}. In the process of working on implementing the model, we extended the existing library with the following functionality:
\begin{itemize}
    \item functionality to filter the table of ngrams (cut off ngram sequences with counts below a certain threshold)
    \item adding writing to and reading ngram models from file
\end{itemize} 
\chapter{Evaluation}
\label{chap:Evaluation}

\epigraph{"It doesn't matter how beautiful your theory is, it doesn't matter how smart you are. If it doesn't agree with experiment, it's wrong."}{--- \textup{Richard Feynman}}

In this chapter, we go over different approaches and challenges for evaluating code completion, and what results we got from our experiments.

\section{Evaluation Overview}
\label{sec:Evaluation-Overview}
Evaluation is an essential part of any research process. In the Richard Feynman quote above, the evaluation \textit{is} the experiment he refers to. To paraphrase and make it more relevant to software engineering, it does not matter what kind of tool you implemented and how well do you think it works, unless there are clear metrics the implementation had been measured on. Now, there are two main types of evaluation you can perform: quantitative and qualitative.

Quantitative, as can be guessed from its name, refers to the type of an evaluation experiment that is based on quantifiable data and uses objective metrics. This is the predominant evaluation approach in the natural sciences and computer science, i.e. measuring physical parameters in the case of the former, or testing performance score in the latter.

Qualitative methods, more native to fields such as sociology or psychology, are now being increasingly used in evaluation of computer systems and software tooling. The main objective is to get the "feel" for the system and deeper explore the human interaction experience, with the data gathered primarily from observations and interviews and subsequently analysed.

Conventionally, a proper qualitative evaluation requires an extensive set-up, involving multiple people, recording their coding process, conducting feedback interviews, and so on. However, due to lack of time, the qualitative evaluation I conducted for this project is more minimal. It involves manually selecting a number of common use cases and interacting with them as a user normally would, while recording the results and their meaning.

\section{Challenges in Evaluating Completion}
\label{sec:Evaluation-Challenges}
\oz{This section does not discuss the challenges of evaluating completion. There is no useful information about quantitative evaluation, and no information about completion and the complexity of telling if it's good or bad.}

\section{Quantitative Evaluation}
\label{sec:Evaluation-Quantitative}
\subsection{Experiment}
The quantitative evaluation approach \badstyle{we decided to take} is based on the paper by \cite{Robb08a}, wherein they used the change history data to evaluate the performance of their completion.

The idea with this approach is to recreate the coding process as it would happen naturally, and test the completion during that. Just as if a developer is typing and invoking a code completion engine at every step, \repetition{our completion gets triggered for every token with the cursor between the 2nd and the 8th position} (this range also taken from the \cite{Robb08a} paper). \repetition{Starting with every second symbol of an alphabetic token, we invoke the completion} and compare whether the result suggested at that step matches the one that followed in real history.

We used a dataset of methods from the Pharo IDE source code \insertion{collected by \cite{Zait20a}} \oz{you must put a reference, otherwise it's impossible for your readers to know what is this dataset, where did it come from, how many methods does it contain, etc.}, and for each of them we repeat this process of calling the completion at each step and performing the following:
\begin{itemize}
    \item check whether the correct completion was present in the list of suggestions
    \item if it was, we only consider the results successful if the correct completion was also present in the top-10 suggestions
    \item for those cases we record the current sequence (i.e. what has been "typed in" up to that point), prefix length, the actual correct match, and its position in the list of suggestions
    \item we also calculate the accuracy (average success rate) for all the methods evaluated (i.e. out of all the attempts to complete the history, how many of them had the expected match in the first 10 results)
\end{itemize}

\subsection{Results}
Table~\ref{table:quan1} presents the accuracy for each of the classes and sorting strategies tested.

\oz{Explain what are the three models that you are comparing: unigram, bigram, and alphabetic. Why did you choose those three, why not 3-gram and 4-gram models? This is the first time when alphabetic sorting should be mentioned (as a "random baseline")}

\begin{table}[H]
    \centering
    \begin{tabular}{lrrrr}
    \hline
    \textbf{Class} & \textbf{\# of Methods} & \textbf{Alphabetic} & \textbf{Unigram} & \textbf{Bigram} \\ \hline
    OrderedCollection & 64 & 0.32 & 0.38 & 0.30 \\ 
    CompletionEvaluation & 9 & 0.24 & 0.25 & 0.22 \\ 
    FrequencyCompletionSorter & 2 & 0.48 & 0.58 & 0.47 \\ 
    NgramModel & 21 & 0.32 & 0.36 & 0.29 \\ \hline
    \end{tabular}
\caption{Quantitative evaluation, average success ratio}
\label{table:quan1}
\end{table}

\oz{How to make your tables cleaner: (1) Round your numbers to 1-3 decimal places; (2) Align text to the left and numbers to the right; (3) Don't use CAPS; (4) Don't use vertical lines; (5) Use less horizontal lines. If you like it, apply those changes to every table}

\oz{And general rules of using tables: (1) Every table must be referenced in text, for example, "In Table~\ref{table:quan1} you can see..."; (2) The text around that reference must be sufficient to understand the table (describe every column, discuss the values); (3) Caption of every table must be sufficient to understand that table without searching for its reference in the main text (captions such as "Qualitative evaluation, example 1" are not enough). Same rules apply to all figures.}

We can see that the unigram sorter performs better than the alphabetic sorter, as expected. However, surprisingly the bigram sorter performs worse than the alphabetic sorter. \oz{And it's not that surprising considering that you are testing it on very small classes.} \notclear{One reason might be the problems with training the model in regards to delimiters and other kinds of non-alphabetic tokens, which are plentiful in the code, and yet are difficult to tackle for correct training.} \oz{It's not clear what you mean.}

\section{Qualitative Evaluation}
\label{sec:Evaluation-Qualitative}
\subsection{Experiment}
Even though the hard numbers seem to show that the bigram strategy is not any good, it really does feel as if it is actually better than the alphabetic approach when used during a regular coding process. To demonstrate this, with no bias towards any of the sorting strategies, we preselected 10 cases for which we were interested to see completion suggestions and how accurate they were, and manually tested each of them. Below are the explanations and results for each of the 10 cases.

\subsection{Results}
Note: if the \textit{completion suggestion} and \textit{position} columns contain '-' for any sorting strategy, it means that in the first 10 suggestions we did not see the results we were expecting to see, based on the context at hand.

It should be taken into consideration, however, that all the results are only recorded after two symbols were typed in, and having correct suggestions with such a short prefix can be considered best-case scenario (as the correct suggestions would definitely be more likely to appear with a longer prefix). Ergo, the absence of correct suggestions at this point does not signal a complete failure of the completion engine or a sorting strategy used, but rather hints that there is more work required to get to the correct suggestions. It is unideal, and is exactly the reason for trying to improve the completion tool.

\begin{description}
\item [Example 1:] \hfill
\begin{lstlisting}
    col := OrderedCollection new.
    col ad
\end{lstlisting}
This was typed into the Playground and the results we were interested in were recorded with the cursor directly after \textit{ad}. For this context, the completion suggestions we were expecting as developers were \textit{add:} and, following, \textit{addAll:}. Here's what we got:
\begin{table}[H]
    \centering
    \begin{tabular}{|l|l|l|}
    \hline
    \textbf{SORTER TYPE} & \textbf{COMPLETION SUGGESTION} & \textbf{POSITION} \\ \hline
    alphabetic & - & - \\ \hline
    unigram & \begin{tabular}[c]{@{}l@{}}add:\\ addAll:\end{tabular} & \begin{tabular}[c]{@{}l@{}}1\\ 2\end{tabular} \\ \hline
    bigram & - & - \\ \hline
    \end{tabular}
\caption{Qualitative evaluation, example 1}
\label{table:qual1}
\end{table}
Both alphabetic and bigram sorters didn't manage to provide any relevant results in the top 10 suggestions, whereas the result from the unigram was exactly what we wanted.

\item [Example 2:] \hfill
\begin{lstlisting}
    something := 1.
    (something = 1) if
\end{lstlisting}
Another example coded in the Playground, with the cursor at \textit{if}. The expected results would be \textit{ifTrue:}, \textit{ifTrue:ifFalse:}, and \textit{ifFalse:}.
\begin{table}[H]
    \centering
    \begin{tabular}{|l|l|l|}
    \hline
    \textbf{SORTER TYPE} & \textbf{COMPLETION SUGGESTION} & \textbf{POSITION} \\ \hline
    alphabetic & - & - \\ \hline
    unigram & \begin{tabular}[c]{@{}l@{}}ifTrue:\\ ifTrue:ifFalse:\\ ifFalse:\end{tabular} & \begin{tabular}[c]{@{}l@{}}1\\ 2\\ 3\end{tabular} \\ \hline
    bigram & ifTrue: & 5 \\ \hline
    \end{tabular}
\caption{Qualitative evaluation, example 2}
\label{table:qual2}
\end{table}
The alphabetic sorter once again did not suggest anything of interest right away, the unigram demonstrated an exemplary performance, and bigram, while less than ideal, still suggested one of the correct options in the top 10 position (albeit a bit lower, which would require either typing more symbols or pressing the down arrow key).

\item [Example 3:] \hfill
\begin{lstlisting}
    TokenProcessing >> returnProcessedData
        | tokensDataFrame |
        tokensDataFrame := self re
\end{lstlisting}
This was typed in the Editor, in the method \textit{returnProcessedData} of class \textit{TokenProcessing}. This is done to try out some more elaborate examples where the local context can play a role. Here, as \textit{rejectInvalidTokens}, \textit{replaceWithPlaceholders} and \textit{readFile} are also methods on the instance side of the \textit{TokenProcessing} class, we expect to get their names as suggestions after typing \textit{self}. Particularly, as the cursor is positioned after \textit{re}, we expect to get only these 3 methods, as they are the ones starting with \textit{re}.
\begin{table}[H]
    \centering
    \begin{tabular}{|l|l|l|}
    \hline
    \textbf{SORTER TYPE} & \textbf{COMPLETION SUGGESTION} & \textbf{POSITION} \\ \hline
    alphabetic & readFile & 3 \\ \hline
    unigram & - & - \\ \hline
    bigram & \begin{tabular}[c]{@{}l@{}}rejectInvalidTokens\\ replaceWithPlaceholders\\ readFile\end{tabular} & \begin{tabular}[c]{@{}l@{}}7\\ 8\\ 9\end{tabular} \\ \hline
    \end{tabular}
\caption{Qualitative evaluation, example 3}
\label{table:qual3}
\end{table}
It seems that here the unigram approach fell short, whereas the bigram approach, trained on prioriting local method name after \textit{self}, has performed quite well. In the alphabetic approach, the other two method suggestions seem to be buried below the first 10, whereas \textit{readFile} gets a better position.

\item [Example 4:] \hfill
\begin{lstlisting}
    TokenProcessing >> replaceWithPlaceholders
        | data tokens tokenTypes arr |
        data := self returnProcessedData.
        to
\end{lstlisting}
Typing \textit{to} we are expecting to have \textit{tokens} and \textit{tokenTypes}, as those are temporary variables we have declared and we might want to initialise them with some value. Here is what we got suggested:
\begin{table}[H]
    \centering
    \begin{tabular}{|l|l|l|}
    \hline
    \textbf{SORTER TYPE} & \textbf{COMPLETION SUGGESTION} & \textbf{POSITION} \\ \hline
    alphabetic & \begin{tabular}[c]{@{}l@{}}tokenTypes\\ tokens\end{tabular} & \begin{tabular}[c]{@{}l@{}}1\\ 2\end{tabular} \\ \hline
    unigram & \begin{tabular}[c]{@{}l@{}}tokens\\ tokenTypes\end{tabular} & \begin{tabular}[c]{@{}l@{}}1\\ 2\end{tabular} \\ \hline
    bigram & \begin{tabular}[c]{@{}l@{}}token\\ tokenTypes\end{tabular} & \begin{tabular}[c]{@{}l@{}}1\\ 2\end{tabular} \\ \hline
    \end{tabular}
\caption{Qualitative evaluation, example 4}
\label{table:qual4}
\end{table}
The order in the unigram and bigram cases is a bit better, as having a shorter word suggested first is preferable, but in general all three methods worked well in this one.

\item [Example 5:] \hfill
\begin{lstlisting}
    UnigramTableCreator ne
\end{lstlisting}
Typing this in the Playground, we expect \textit{new} and maybe even \textit{new:}, depending on the context.
\begin{table}[H]
    \centering
    \begin{tabular}{|l|l|l|}
    \hline
    \textbf{SORTER TYPE} & \textbf{COMPLETION SUGGESTION} & \textbf{POSITION} \\ \hline
    alphabetic & \begin{tabular}[c]{@{}l@{}}new\\ new:\end{tabular} & \begin{tabular}[c]{@{}l@{}}2\\ 3\end{tabular} \\ \hline
    unigram & \begin{tabular}[c]{@{}l@{}}new\\ new:\end{tabular} & \begin{tabular}[c]{@{}l@{}}1\\ 2\end{tabular} \\ \hline
    bigram & \begin{tabular}[c]{@{}l@{}}new:\\ new\end{tabular} & \begin{tabular}[c]{@{}l@{}}4\\ 8\end{tabular} \\ \hline
    \end{tabular}
\caption{Qualitative evaluation, example 5}
\label{table:qual5}
\end{table}
The best positions are for the unigram sorter, but in general all the sorters suggested the desired options in the top 10.

\item [Example 6:] \hfill
\begin{lstlisting}
    in
\end{lstlisting}
In the Pharo IDE it is also possible to get completion suggestions even when creating a new method, and in this case we went to the instance side of a class and started typing \textit{in} as method name, wanting to get \textit{initialize}, as it is a very common method one might want to create in a class.
\begin{table}[H]
    \centering
    \begin{tabular}{|l|l|l|}
    \hline
    \textbf{SORTER TYPE} & \textbf{COMPLETION SUGGESTION} & \textbf{POSITION} \\ \hline
    alphabetic & - & - \\ \hline
    unigram & initialize & 2 \\ \hline
    bigram & - & - \\ \hline
    \end{tabular}
\caption{Qualitative evaluation, example 6}
\label{table:qual6}
\end{table}
The unigram sorter is the only that was able to suggest \textit{initialize} from only 2 characters typed in, and at the 2 position!

\item [Example 7:] \hfill
\begin{lstlisting}
    ExperimentalNgramTest >> testMostLikelyContinuations
    | expected actual word |
    word := 'lorem'.
    actual := model mostLikelyContinuations: word asNgram top: 3.
    expected := #('ipsum' 'dolor' 'lorem').
    self as
\end{lstlisting}
In test methods having an \textit{assert:equals:} call is quite common, and as there are a lot of tests being written, we are interested in having a completion suggestion that is as accurate as possible after \textit{as}.
\begin{table}[H]
    \centering
    \begin{tabular}{|l|l|l|}
    \hline
    \textbf{SORTER TYPE} & \textbf{COMPLETION SUGGESTION} & \textbf{POSITION} \\ \hline
    alphabetic & - & - \\ \hline
    unigram & \begin{tabular}[c]{@{}l@{}}assert:\\ assert:equals:\end{tabular} & \begin{tabular}[c]{@{}l@{}}1\\ 2\end{tabular} \\ \hline
    bigram & assert:equals: & 1 \\ \hline
    \end{tabular}
\caption{Qualitative evaluation, example 7}
\label{table:qual7}
\end{table}
Both the unigram and bigram sorters showed good results.

\item [Example 8:] \hfill
\begin{lstlisting}
    CompletionEvaluation >> averageSuccessRatio
    | sumS sumA |
	sumS := 0.
	sumA := 0.
	countSuccess do: [ :each | su ]
\end{lstlisting}
Here we are interested in completion inside the block, so our cursor is after \textit{su} before the bracket \textit{]}. Since inside the block we want to do assignment for the temporary variables, we are expecting \textit{sumA} and \textit{sumS} as suggestions.
\begin{table}[H]
    \centering
    \begin{tabular}{|l|l|l|}
    \hline
    \textbf{SORTER TYPE} & \textbf{COMPLETION SUGGESTION} & \textbf{POSITION} \\ \hline
    alphabetic & \begin{tabular}[c]{@{}l@{}}sumA\\ sumS\\ super\end{tabular} & \begin{tabular}[c]{@{}l@{}}1\\ 2\\ 3\end{tabular} \\ \hline
    unigram & \begin{tabular}[c]{@{}l@{}}super\\ sumA\\ sumS\end{tabular} & \begin{tabular}[c]{@{}l@{}}1\\ 2\\ 3\end{tabular} \\ \hline
    bigram & \begin{tabular}[c]{@{}l@{}}super\\ sumA\\ sumS\end{tabular} & \begin{tabular}[c]{@{}l@{}}1\\ 2\\ 3\end{tabular} \\ \hline
    \end{tabular}
\caption{Qualitative evaluation, example 8}
\label{table:qual8}
\end{table}

\item [Example 9:] \hfill
\begin{lstlisting}
    RBScanner sc
\end{lstlisting}
In the Playground we want to get suggested method names on the class side of \textit{RBScanner} that begin with \textit{sc}. There are two: \textit{scanTokens:} and \textit{scanTokenObjects:}.
\begin{table}[H]
    \centering
    \begin{tabular}{|l|l|l|}
    \hline
    \textbf{SORTER TYPE} & \textbf{COMPLETION SUGGESTION} & \textbf{POSITION} \\ \hline
    alphabetic & \begin{tabular}[c]{@{}l@{}}scanTokenObjects:\\ scanTokens:\end{tabular} & \begin{tabular}[c]{@{}l@{}}1\\ 2\end{tabular} \\ \hline
    unigram & \begin{tabular}[c]{@{}l@{}}scanTokenObjects:\\ scanTokens:\end{tabular} & \begin{tabular}[c]{@{}l@{}}1\\ 2\end{tabular} \\ \hline
    bigram & \begin{tabular}[c]{@{}l@{}}scanTokens:\\ scanTokenObjects:\end{tabular} & \begin{tabular}[c]{@{}l@{}}1\\ 2\end{tabular} \\ \hline
    \end{tabular}
\caption{Qualitative evaluation, example 9}
\label{table:qual9}
\end{table}
On the one hand, all sorters showed essentially the same result; however, on the other hand, \textit{scanTokens:} is more likely to be used next to \textit{RBscanner} and it is a shorter token, so the bigram result is slightly superior.

\item [Example 10:] \hfill
\begin{lstlisting}
    UnigramTableCreator >> writeToFile
    | data stream |
    data := self frequencies.
    stream := TokenProcessing frequenciesFile writeStream.
    (NeoCSVWriter on: stream)
        nextPut: #(key value);
        ne
\end{lstlisting}
Now, for an even more creative example, in Pharo there can be multiple messages sent to a receiver, as long as they are all separated by a semicolon \textit{;}. Hence, having a cursor after \textit{ne}, we know we are sending a message to \textit{NeoCSVWriter on: stream} and want to get suitable suggestions, which are \textit{nextPut:} and \textit{nextPutAll:}.

In the table of results below we see that the unigram sorter got the best result.
\begin{table}[H]
    \centering
    \begin{tabular}{|l|l|l|}
    \hline
    \textbf{SORTER TYPE} & \textbf{COMPLETION SUGGESTION} & \textbf{POSITION} \\ \hline
    alphabetic & - & - \\ \hline
    unigram & \begin{tabular}[c]{@{}l@{}}nextPutAll:\\ nextPut:\end{tabular} & \begin{tabular}[c]{@{}l@{}}2\\ 3\end{tabular} \\ \hline
    bigram & - & - \\ \hline
    \end{tabular}
\caption{Qualitative evaluation, example 10}
\label{table:qual10}
\end{table}
\end{description}

To briefly sum up the results of the manual qualitative evaluation, the unigram sorter is still the best one out of the three, the same as in the quantitative evaluation. However, as seen in these examples, the bigram sorter is not very far behind, also more often than not prioriting the desired completions, which disagrees with the quantitative evaluation. So, while based on cold data, the bigram sorter is the worst of the three \textit{on average}, from the manual evaluation it seems that it gives perfectly adequate results even for more 'tricky' examples.

\section{Summary}
\label{sec:Evaluation-Summary}
\begin{itemize}
    \item Quantitative (automatic) evaluation involves evaluating the model on \badstyle{some benchmarks}, whereas the qualitative (manual) approach involves evaluating the results \badstyle{on the go}, as the program is being normally executed.
    \item \remove{There are some challenges when properly setting up the qualitative approach, as it involves a lot of resources; we didn't have them and so conducted a minimalist version of the approach.} \oz{It's not true. We had them. You were doing internship in a team with over 20 experienced researchers and developers. And a highly responsive Pharo community with hundreds of developers willing to participate in experiments for free. So it is unfair to say that we don't have resources. You didn't have time to use those resources. So just remove this statement}
    \item Based on the evaluation tool that we created and used to evaluate average success rate (i.e. out of all the attempts to complete the history, how many of them had the expected match in the first 10 results) for each sorting strategy, the unigram sorter performed better than the alphabetic sorter but, surprisingly, the bigram one performed worse.
    \item For the qualitative evaluation, we picked 10 interesting completion cases (with no bias towards any strategy) and tested them in the real environment and recorded the positions of completion suggestions we were interested in. Based on this, the unigram approach was once again the best, but this time the bigram approach \badstyle{followed not far behind}, and the alphabetic was worse.
    \item \remove{Hence we can officially state that} \remove{the unigram sorter is a clear improvement on the alphabetic approach, whereas the bigram sorter seems to lack the accuracy.} \oz{Move this to the Conclusion. But don't say "clear improvement". How do you define "clear"? For me it's not clear at all}
\end{itemize}
\chapter{Conclusion}
\label{chap:Conclusion}
In this work, I proposed a machine learning-based technique for improving the completion engine in the Pharo IDE. The exact implementation uses the n-gram language models to sort the completion candidates that are suggested to the user as they type. Specifically, I implemented two sorting strategies: one based on the unigram model and another on the bigram model.

Based on the evaluations conducted and the actual usage of the implemented sorters in the Pharo IDE, one of the implementations, the unigram sorter, has been shown to be: (1) fast, which means that it can be comfortably used by Pharo developers when coding, and (2) effective, which means it gives significantly better, or more relevant, results than the default Pharo IDE sorter.

Therefore, we can conclude that the main goal, improving the Pharo code completion with the help of an n-gram based sorting implementation, has been achieved.

\section{Discoveries}
\label{sec:Conclusion-Discoveries}
Through completing this work, I am now able to answer the research questions initially stated in Section \ref{sec:Introduction-Approach}:
\begin{RQ}
    \item \textbf{Can we improve the accuracy of code completion in the Pharo IDE by sorting candidate completions with n-gram language models?} As a result of the evaluation performed, the unigram based sorter has been shown to have a significantly better result than the default sorter in the Pharo IDE. The implementation is also fast enough for comfortable developer usage and is available by loading the CompletionSorting library available at \url{https://github.com/myroslavarm/CompletionSorting}.
    \item \textbf{How can we effectively evaluate the results of code completion enhanced by different sorting strategies?} For the quantitative evaluation, inspired by \cite{Robb08a}, I simulated the completion process as it would happen naturally, by generating source code sequences at various stages of typing, and comparing the results proposed to the ones in the codebase.
    
    The qualitative approach allowed me to experimentally test the suggestions each sorting strategy proposes by using the tool as it would be normally used by a developer, and compare the results first-hand.
\end{RQ}

\section{Directions of Future Work}
\label{sec:Conclusion-FutureWork}
\subsection{Enhancing the Bigram Sorter}
As can be seen in Chapter \ref{chap:Evaluation}, the bigram sorter did not perform as well as expected. Contrary to the intuition that the higher order of n-gram should work better, the bigram performed much worse than the unigram. Additionally, it also seemed to give less relevant suggestions than the alphabetic sorter.

Hence, for future work, it would be useful to see what exactly went wrong with the implementation. It could be that there is a mismatch between how source code is parsed and tokenised in the training and test data. Or it could be a matter of managing the delimiters incorrectly, as punctuation in source code has a significant influence on the contextual meaning of any part of code. It could also be an issue with the NgramModel library, which was only used for training the bigram model, whereas the unigram model was implemented using a different approach.

In any case, this would be a good idea for the continuation of this project: analyse the issues of the existing bigram implementation and try to fix them and improve the performance of the bigram-based sorter. 
\subsection{Conducting a More Extensive Evaluation}
For the quantitative evaluation, it would be useful to evaluate the results more precisely. Meaning that instead of an average accuracy, perhaps a more sophisticated formula could be used. For instance, the idea by \cite{Robb08a} involves making a note of the exact positions and prefix lengths and ranking completions, which give more relevant results for shorter prefixes, higher.

For the qualitative approach, a more extensive case study with multiple participants testing the actual completion in the IDE and reporting their feedback would be a good next step towards a more thorough evaluation of the results.


%----------------------------------------------------------------------------------------
%	THESIS CONTENT - APPENDICES
%----------------------------------------------------------------------------------------

\appendix % Cue to tell LaTeX that the following "chapters" are Appendices

% Include the appendices of the thesis as separate files from the Appendices folder
% Uncomment the lines as you write the Appendices

%\include{Appendices/AppendixA}
%\include{Appendices/AppendixB}
%\include{Appendices/AppendixC}

%----------------------------------------------------------------------------------------
%	BIBLIOGRAPHY
%----------------------------------------------------------------------------------------

\printbibliography[heading=bibintoc]

%----------------------------------------------------------------------------------------

\end{document}  
